\documentclass{report} 
\title{Strogatz}
\date{Started 20th November 2025}
\author{Malcolm}
\usepackage{amsmath} %import math
\usepackage{mathtools} %more math
\usepackage{amssymb} %for QED symbol
\usepackage{amsthm} %
\usepackage{mathrsfs} %
\usepackage{bm}%bold math
\usepackage{graphicx} %import imaging

\graphicspath{{./images/}} %set imaging path
\newtheorem*{theorem}{Theorem}
\theoremstyle{definition}
\newtheorem*{definition}{Definition}


\begin{document}
\maketitle

\tableofcontents

\chapter{Flows on a line}
\section{Introduction}
We have the general system as
\begin{align*}
\dot x_1=&f_1(x_1,\ldots,x_n)\\
&\vdots\\
\dot x_n=&f_n(x_1,\ldots,x_n)
\end{align*}
its solutions can be intuitively visualised as trajectories flowing through an $n$-dimensional phase space. In these early sections we start with the simple case where $n=1$; then we get a single
equation of the form
\begin{equation*}
\dot x=f(x)
\end{equation*}
We call such systems one-dimensional or first-order systems.\\
\vspace{1mm}\\
We clarify the following:
\begin{enumerate}
\item The word \textit{system} is being used here in the context of a dynamical system, not in the classical sense of a collection of two or more equations. Thus a single equation can be called a
`system'.
\item We do not allow $f$ to depend explicitly on time. Time-dependent or `nonautonomous' equations of the form $\dot x=f(x,t)$ are more complicated because they require two variables. Such a system
in this case would be regarded as a two-dimensional system, to be discussed later.
\end{enumerate}
\newpage

\section{Geometric intuition}
\subsection{Instructive example}
Consider the following nonlinear differential equation:
\begin{equation*}
\dot x=\sin x
\end{equation*}
This particular equation can be solved in closed form; using separation of variables then integration we get
\begin{equation*}
t=\ln\left|\frac{\csc x_0+\cot x_0}{\csc x+\cot x}\right|
\end{equation*}
where $x=x_0$ at $t=0$. Although the result is exact, it isn't really possible to interpret. For instance consider the following
\begin{enumerate}
\item Suppose $x_0=\pi/4$. Describe the qualitative features of the solution $x(t)$ for all $t>0$. What happens as $t\to\infty$?
\item For an arbitrary initial condition $x_0$, what is the behaviour of $x(t)$ as $t\to\infty$?
\end{enumerate}
The formula above is not transparent. We consider instead a graphical analysis; the differential equation $\dot x=\sin x$ represents a \textit{vector field} on the line: it dictates the velocity
vector $\dot x$ at each $x$. To sketch this vector field, it is convenient to plot $\dot x$ vs $x$, drawing arrows on the $x$-axis to indicate the corresponding velocity vector at each $x$.
The arrows point to the right when $\dot x>0$ and to the left when $\dot x<0$:
\begin{center}
\includegraphics[width=10cm]{1}
\end{center}
The \textit{flow} is to the right when $\dot x>0$ and to the left when $\dot x<0$; at points where $\dot x=0$, there is no flow, and such points are therefore called \textit{fixed points}. Solid
dots represent \textit{stable} fixed points (often called \textit{attractors/sinks}, bscause the flow is toward them). Open dots represent \textit{unstable} fixed points (also known as \textit{repellers/sources}).\\
(next page)\newpage
\noindent\textbf{Cont.}\\
We had \begin{center}
\includegraphics[width=9cm]{1}
\end{center}
We start at $x_0$ and watch how $x$ is carried along by the flow. Consider now the questions proposed earlier.
\begin{enumerate}
\item Starting at $x_0=\pi/4$, the flow moves to the right faster and faster until $x$ crosses $\pi/2$ (where $\sin x$ reaches its maximum). Then the flow slows down and eventually approaches
the stable fixed point $x=\pi$ from the left. As such the qualitative form of the solution is as follows
\begin{center}
\includegraphics[width=7cm]{2}
\end{center}
See that the curve is concave up at first, then concave down; this change in concavity corresponds to the initial acceleration for $x<\pi/2$, followed by the deceleration toward $x=\pi$.\\
(next page)\newpage
\item The same reasoning applies to any initial condition $x_0$. If $\dot x>0$ initially, the flow is to the right and $x$ asymptotically approaches the nearest stable point, and similarly to the left for the
other case where $\dot x<0$. If $\dot x=0$, then $x$ remains constant. The qualitative form of the solution for any initial condition can be sketched: 
\begin{center}
\includegraphics[width=6cm]{3}
\end{center}
\end{enumerate}
\noindent Note that the picture can't tell us certain \textit{quantitative} things: for instance we don't know the time at which the speed $|\dot x|$ is greatest. But in many cases the \textit{qualitative}
information is what we care about. In this case these sketches are useful.
\newpage

\subsection{Fixed points and stability}
The ideas discussed can be extended to any one-dimensional system $\dot x=f(x)$. We draw the graph of $f(x)$.
\begin{center}
\includegraphics[width=8cm]{4}
\end{center}
As before, the flow is to the right where $f(x)>0$ and to the left when $f(x)<0$. To find the solution to $\dot x=f(x)$ starting from an arbitrary initial condition $x_0$, we start at $x=x_0$ and 
watch how it is carried along by the flow.\\
\vspace{1mm}\\
The point moves along the $x$ axis according to some function $x(t)$. This function is called the \textit{trajectory} based at $x_0$, and it represents the solution of the differential equation
starting from the initial condition $x_0$. A figure like the one above, which shows all the different trajectories of the system, is called a \textit{phase portrait}.\\
\vspace{1mm}\\
The appearance of the phase portrait is controlled by the fixed points $x^*$, defined by $f(x^*)=0$; they correspond to stagnation points of flow. As shown earlier the solid dots represent stable
fixed points while open dots represent unstable fixed points. In terms of trajectories, fixed points represent \textit{equilibrium} solutions (sometimes called steady/constant/rest solutions), 
since if $x=x^*$ initially it stays there for all of time.\\
\vspace{1mm}\\
An equilibrium is defined to be stable if all sufficiently small disturbances away from it damp out in time, so stable equilibria are 
represented geometrically by stable fixed points. Conversely, unstable equilibria, in which disturbances grow over time, are represented by unstable points.\\
(next page)\newpage
\noindent\textbf{Instructive examples}\\
Say we want to find all the fixed points for $\dot x=x^2-1$, and then classify their stability.\\
\vspace{1mm}\\
Here $f(x)=x^2-1$. To find the fixed points we set $f(x^*)=0$ and solve for $x^*$; thus $x^*=\pm1$. To determine stability we plot $x^2-1$ and then sketch the vector field:
\begin{center}
\includegraphics[width=7cm]{5}
\end{center}
The flow is to the right when $x^2-1>0$ and to the left where $x^2-1<0$. $x^*=-1$ is stable and $x^*=1$ is unstable.\\
\vspace{1mm}\\
Note that the definition of stable equilibrium is based on \textit{small} disturbances; certain large disturbances may fail to decay. In this case see that all small disturbances to 
$x^*=-1$ will decay, but a large disturbance that sends $x$ to the right of $x=1$ will \textit{not} decay---and will in fact be repelled out to $+\infty$. To emphasise this we sometimes say that
$x^*=-1$ is \textit{locally stable}, but not globally stable.\\
\vspace{1mm}\\
On the other hand consider the vector field for the RC circuit equation
\begin{equation*}
\dot Q=f(Q)=\frac{V_0}{R}-\frac{Q}{RC}
\end{equation*}
\begin{center}
\includegraphics[width=5cm]{6}
\end{center}
In this case the flow is always towards $Q^*$; it is a \textit{stable} fixed point. In fact, it is \textit{globally stable}, in the sense that it is approached from \textit{all} initial 
conditions.\\
(next page)\newpage
\noindent\textbf{Another example}\\
Say we want to sketch the phase portrait corresponding to $\dot x=x-\cos x$, and determine the stability of all the fixed points.\\
\vspace{1mm}\\
Our approach would be to plot $x-\cos x$ and then sketch the associated vector field. This method is valid, but it requires one to figure out what the graph of $x-\cos x$ looks like. 
An easier approach would be to plot both $y=x$ and $y=\cos x$ separately; we plot both graphs on the same axes and then observe that they intersect at exactly one point:
\begin{center}
\includegraphics[width=8cm]{7}
\end{center}
This intersection corresponds to a fixed point since $x^*=\cos x^*$ and therefore $f(x^*)=0$. When the line lies above the cosine curve, we have $x>\cos x$ and so $\dot x>0$; the flow is to the
right. Similarly, the flow is to the left when the line is below the cosine curve.\\
\vspace{1mm}\\
Hence $x^*$ is the only fixed point and is unstable. See that we can classify its stability even though we don't know exactly where it is.
\newpage

\section{Linear stability analysis}
One may want to have a more quantitative measure of stability, such as the rate of decay to a stable fixed point. This sort of information
may be obtained by \textit{linearising} about a fixed point.\\
\vspace{1mm}\\
Let $x^*$ be a fixed point, and let
\begin{equation*}
x(t)=x^*+u(t)
\end{equation*}
where 
\begin{equation*}
u(t)=x(t)-x^*
\end{equation*}
differentiation yields
\begin{equation*}
\dot u=\frac{d}{dt}(x-x^*)=\dot x
\end{equation*}
Thus $\dot u=\dot x=f(x)=f(x^*+u)$. Now using Taylor's expansion we obtain
\begin{equation*}
f(x^*+u)=f(x^*)+f'(x^*)u+O(u^2)
\end{equation*}
where $O(u^2)$ denotes quadratically small terms in $u$. Finally observe that $f'(x^*)=0$ since $x^*$ is a fixed point. Hence
\begin{equation*}
\dot u=f'(x^*)u+O(u^2)
\end{equation*}
If $f'(x^*)\neq0$, the $O(u^2)$ terms are negligible in comparison to the linear term, so we may write the approximation
\begin{equation*}
\dot u\approx f'(x^*)u
\end{equation*}
This linear differential equation for $u$ is called the \textit{linearisation} about $x^*$. The linearisation shows that $u(t)$ grows exponentially if $f'(x^*)>0$ and decays exponentially if 
$f'(x^*)<0$. If $f'(x^*)=0$, the $O(u^2)$ can't be ignored and a nonlinear analysis is needed to determine stability.\\
\vspace{1mm}\\
The slope $f'(x^*)$ at the fixed point determines its stability. The importance of the \textit{sign} of $f'(x^*)$ was clear from our graphical approach; now we also
have a measure of \textit{how} stable a fixed point is---determined by the \textit{magnitude} of $f'(x^*)$. Its reciprocal $1/|f'(x^*)|$ is a \textit{characteristic time scale}; it determines the 
time required for $x(t)$ to vary significantly in the neighbourhood of $x^*$.
\newpage

\section{Existence and uniqueness}
\textbf{Instructive example}\\
We can show that the solution to $\dot x=x^{1/3}$ starting from $x_0=0$ is \textit{not} unique.\\
\vspace{1mm}\\
The point $x=0$ is a fixed point, so one obvious solution is $x(t)=0$ for all $t$. 
The suprising fact is that there is another solution. By separating variables and integrating:
\begin{equation*}
\int x^{-1/3}dx=\int dt
\end{equation*}
so $3/2x^{2/3}=t+C$. Imposing an initial condition $x(0)=0$ yields $C=0$, and hence $x(t)=(2/3t)^{3/2}$ is also a solution.\\
\vspace{1mm}\\
To find the source of the non-uniqueness, a hint comes from the vector field:
\begin{center}
\includegraphics[width=5cm]{8}
\end{center}
See that the fixed point $x^*=0$ is very unstable---the slope of $f(0)$ is infinite---a jump discontinuity in $x$. We state the existence and uniqueness theorem as follows\\
\vspace{1mm}\\
\textbf{Existence and uniqueness Theorem}: Consider the initial value problem
\begin{equation*}
\dot x=f(x),\quad x(0)=x_0
\end{equation*}
Suppose that $f(x)$ and $f'(x)$ are continuous on an open interval $R$ of the $x$ axis, and suppose that $x_0$ is a point in $R$. Then the initial
value problem has a solution $x(t)$ on some time interval $(-\tau,\tau)$ about $t=0$, and the solution is unique.\\
\vspace{1mm}\\
This theorem says that if $f(x)$ is \textit{smooth enough}, then solutions exist and they are unique. Even so, note that there's \textit{no guarantee that solutions exist for all time}.\\
(next page)\newpage
\noindent\textbf{Instructive example}\\
Consider the initial value problem $\dot x=1+x^2$, $x(0)=x_0$. We have $f(x)=1+x^2$. This function is continuous and has a continuous derivative for all of $x$. Hence the theorem tells 
us that solutions exist and are unique for any initial condition $x_0$. 
But \textit{the theorem does not say that the solutions exist for all time}; they are only guaranteed to exist in a (possibly very short) time interval around $t=0$.\\
\vspace{1mm}\\
For instance consider the case where $x(0)=0$. The problem can can be solved analytically by separation of variables to yield
\begin{equation*}
\tan^{-1}x=t+C
\end{equation*}
where the initial condition would imply $C=0$, so $x(t)=\tan t$ is the solution. But notice that this solution only exists for $-\pi/2<t<\pi/2$, since $x\to\pm\infty$ as $t\to\pm\pi/2$. 
Outside this interval there is no solution to the initial value problem for $x_0=0$ (solutions must be continuous).\\
\vspace{1mm}\\
This system has solutions that reach infinity in \textit{finite} time. This phenomenon is called \textit{blow-up}.

\section{Impossibility of Oscillations}
In all examples so far, all trajectories either approached a fixed point, or diverged to $\pm\infty$. In fact, those are the \textit{only} things that can happen to a vector field on the real line.
Trajectories are forced to increase or decrease \textit{monotonically}; overshoot and damped oscillations can never occur in a first-order system. Undamped oscillations are also impossible.
There are no periodic soluutions to $\dot x=f(x)$. These general results are fundamentally topological in origin. 
\newpage

\section{Potentials}
For another way to visualise the dynamics of the first order system $\dot x=f(x)$, we picture a particle sliding down the walls of a potential well, where the \textit{potential} $V(x)$ is defined
by
\begin{equation*}
f(x)=-\frac{dV}{dx}
\end{equation*}
The negative sign in the definition of $V$ implies that the particle always moves ``downhill'' as the motion proceeds.
\begin{center}
\includegraphics[width=8cm]{9}
\end{center}
To see this, we can use the chain rule to obtain
\begin{equation*}
\frac{dV}{dt}=\frac{dV}{dx}\frac{dx}{dt}
\end{equation*}
Hence by the definition of the potential
\begin{equation*}
\frac{dV}{dt}=-\left(\frac{dV}{dx}\right)^2\leq0
\end{equation*}
Thus $V(t)$ \textit{decreases along trajectories} (in the sense that it is strictly decreasing with time), and so the particle always moves toward lower potential. 
Of course, if the particle happens to be at an \textit{equilibrium} point where $dV/dx=0$, then $V$ remains constant. This is to be expected, since $dV/dx=0$ implies $\dot x=0$.\\
\vspace{1mm}\\
In this sense, local minima of $V(x)$ correspond to \textit{stable} fixed points, as we'd expect intuitively, and local maxima correspond to unstable fixed points.\\
(next page)\newpage
\noindent\textbf{Example}\\
Consider the system $\dot x=-x$. We have $-dV/dx=-x$. The general solution is $V(x)=\frac{1}{2}x^2+C$. Choosing $C=0$ for convenience, the graph of $V(x)$ is
\begin{center}
\includegraphics[width=4.5cm]{10}
\end{center}
The only equilibrium point occurs at $x=0$, and it's stable.\\
\vspace{1mm}\\
For another example consider $\dot x=x-x^3$. Solving $-dV/dx=x-x^3$ yields $V=-\frac{1}{2}x^2+\frac{1}{4}x^4+C$. Once again we set $C=0$:
\begin{center}
\includegraphics[width=4.5cm]{11}
\end{center}
The local minima at $x=\pm1$ correspond to stable equilibria, and the local maximum at $x=0$ corresponds to an unstable equilibrium. This is often called a \textit{double-well potential},
and the system is said to be \textit{bistable}, since it has two stable equilibria.
\newpage

\section{Numerical methods}
The Euler's method update rule is
\begin{equation*}
x_{n+1}=x_n+f(x_n)\Delta t
\end{equation*}
where $\Delta t$ is the timestep. Euler's method is not recommended in practice since its approximation is pretty bad. One problem with it is that it estimates the derivative at only the left end
of the time interval between $t_n$ and $t_{n+1}$. A more sensible approach would be to use the 
\textit{average} derivative across this interval.\\
\vspace{1mm}\\
The \textit{improved Euler method} first takes a trial step across the interval as in the generic Euler method. This produces a trial value $\tilde x_{n+1}=x_n+f(x_n)\Delta t$;
the tilde above the $x$ indicates that this is a tentative step, used only as a probe. Now that we've estimated the derivative on both ends of the interval, we average
$f(x_n)$ and $f(\tilde x_{n+1})$, and use that to take the \textit{real} step across the interval.
The improved Euler method is
\begin{align*}
&\tilde x_{n+1}=x_n+f(x_n)\Delta t&\text{(the trial step)}\\
&x_{n+1}=x_n+\frac{1}{2}[f(x_n)+f(\tilde x_{n+1})]\Delta t&\text{(the real step)}
\end{align*}
This method is more accurate in the sense that it tends to make a smaller error $E=|x(t_n)-x_n|$ for a given stepsize $\Delta t$. In both method the error $E\to0$ as $\Delta t\to0$, but the error
decreases \textit{faster} for the improved Euler. The Euler method is considered first order, whereas the improved Euler is second order.\\
\vspace{1mm}\\
Higher order methods have been concocted, but are not necessarily superior, since they are more computationally expensive. In practice, a good balance is achieved by the 
\textit{fourth-order Runge-Kutta method}. We evaluate the following 
\begin{align*}
&k_1=f(x_n)\Delta t\\
&k_2=f(x_n+\frac{1}{2}k_1)\Delta t\\
&k_3=f(x_n+\frac{1}{2}k_2)\Delta t\\
&k_4=f(x_n+k_3)\Delta t
\end{align*}
Then $x_{n+1}$ is given by
\begin{equation*}
x_{n+1}=x_n+\frac{1}{6}(k_1+2k_2+2k_3+k_4)
\end{equation*}
This method generally gives accurate results without requiring an excessively small stepsize $\Delta t$.
\newpage

\chapter{Bifurcations}
The dynamics of vector fields on the line is very limited: all solutions either settle down to equilibrium or head out to $\pm\infty$. Such dynamics are fairly trivial. 
The interest here lies in \textit{dependence on parameters}. The qualitative structure of the flow
can change as parameters are varied. In particular, fixed points can be created or destroyed, or 
their stability can change. These qualitative changes in the dynamics are called \textit{bifurcations}, and the parameter values at which they occur are called \textit{bifurcation points}.

\section{Saddle-Node Bifurcation}
The saddle-node bifurcation is the basic mechanism by which fixed points are \textit{created and destroyed}. As a parameter is varied, two fixed points move toward each other, collide, and mutually
annihilate. The prototypical or \textit{normal form} of a saddle-node bifurcation is given by the first order system
\begin{equation*}
\dot x=r+x^2
\end{equation*}
where $r$ is a parameter which may be positive, negative or zero:
\begin{center}
\includegraphics[width=12cm]{12}
\end{center}
(next page)\newpage
\noindent We have
\begin{equation*}
\dot x=r+x^2
\end{equation*}
where $r$ is a parameter which may be positive, negative or zero:
\begin{center}
\includegraphics[width=12cm]{12}
\end{center}
When $r$ is negative, there are two fixed points, one stable and one unstable. As $r$ approaches 0 from below, the parabola moves up and the two fixed points move toward each other.
When $r=0$ the fixed points coalesce into a half-stable fixed point at $x^*=0$. This type of fixed
point is extremely delicate---vanishing as soon as $r>0$. After which there is no more fixed point.
We say that a \textit{bifurcation} occurred at $r=0$ since the vector fields for $r<0$ and $r>0$ are qualitatively different.\\
\vspace{1mm}\\
\textbf{Illustration}\\
There are several ways to depict such a bifurcation. We can show a stack of vector fields for discrete values of $r$:
\begin{center}
\includegraphics[width=12cm]{13}
\end{center}
(next page)\newpage
\noindent Alternatively we can let $r$ be continuous.
\begin{center}
\includegraphics[width=12cm]{14}
\end{center}
The curve shown is $r=-x^2$; i.e $\dot x=0$, giving the fixed points for different $r$. To distinguish between stable and unstable fixed points, we use a solid line for stable points and a broken
line for unstable ones.\\
\vspace{1mm}\\
However, the most common way to depict the bifurcation is to invert the axes. The idea here is that $r$ plays the role of an independent variable, and so should be plotted horizontally.
\begin{center}
\includegraphics[width=12cm]{15}
\end{center}
Arrows are sometimes included, but not always. This is the \textit{bifurcation diagram} for the saddle-node bifurcation.












\end{document}
